% !TeX program = xelatex
% UTF-8. Standalone: run xelatex "allergic rhinitis.tex" in this directory.
% Inclusion: define \AllergicRhinitisInNotebook before \input in a XePersian document.
% Source: Google Drive Notebook/allergic rhinitis, ascending original filename order.
% pic01.jpg = 20261012_174728683.jpg; 1UWvvunrlol2BDllTc7tXSwr3hFYJk3xF
% pic02.jpg = 20261012_174733697.jpg; 1tN04Yz6Duy6MEhcEBW_KpUhMVedJr-67
% pic03.jpg = 20261012_174737021.jpg; 1PRBROUplLNht5UudHy7kkrmO8-r-gJGO
% pic04.jpg = 20261012_174741948.jpg; 16fUHRWqQ1mqpsI6fzSbtcIyt0GnQSTdq
% pic05.jpg = 20261012_174746233.jpg; 1zIyxUjajWK0vhNuuFlQqSl-K9kXNowVv
% pic06.jpg = 20261012_174750500.jpg; 1pCa0sOH4Js1AFCS9HGTjWpBRB7xfALXr
% pic07.jpg = 20261012_174754320.jpg; 1OBGEi_2po_KetAK2e7w0Btf_9_wilItP
% pic08.jpg = 20261012_174758690.jpg; 1JQRxSGRJ6f2QwZczjzWiC0g65hoUOYUD
% pic09.jpg = 20261012_174802404.jpg; 1bgovOosWvuvWcqPqOOCc6Sz2x_YK5C2f
\ifdefined\AllergicRhinitisInNotebook
  \let\ARfinish\relax
\else
  \def\ARfinish{\end{document}}
  \documentclass[12pt,a4paper]{article}
  \usepackage[margin=22mm]{geometry}
  \usepackage{amsmath,amssymb}
  \usepackage{graphicx}
  \usepackage{xepersian}
  \IfFontExistsTF{Amiri}{\settextfont{Amiri}}{\settextfont{DejaVu Sans}}
  \setlatintextfont{DejaVu Serif}
  \setlength{\parindent}{0pt}
  \setlength{\parskip}{5pt}
  \setlength{\emergencystretch}{3em}
  \begin{document}
\fi
\begingroup
\providecommand{\ARunclear}[1]{\textbf{[خوانش نامطمئن: #1]}}
\providecommand{\ARrange}[2]{\lr{#1--#2}}
\providecommand{\ARpage}[1]{\clearpage\section*{تصویر #1}\noindent
  \textbf{سربرگ و حاشیه:} ۱۸۴ ـ؛ \lr{Subject:} (خالی)؛ \lr{Date: / /}؛ شمارهٔ صفحه: #1؛ نشان پایین برگه: \lr{YASHA}.\par}

\begin{center}
{\Large\bfseries رینیت آلرژیک ـ ۱۸۴}\par
\lr{Nelson -- 22}
\end{center}
\textbf{یادداشت رونویسی:} متن فارسی و انگلیسی، اختصارات و اعداد به ترتیب تصاویر حفظ شده‌اند؛ فلش‌ها به صورت فهرست و رابطه بازنویسی شده‌اند. تأکیدهای رنگی با حروف پررنگ مشخص شده‌اند. مطالب پزشکی منبع اصلاح نشده‌اند. متن خط‌خوردهٔ ناخوانا مشخص شده و تصاویر کامل منبع در پیوست آمده‌اند.

\section*{تصویر ۱}
\textbf{سربرگ و حاشیه:} رینیت آلرژیک ـ ۱۸۴؛ \lr{Subject: Nelson -- 22}؛ \lr{Date: / /}؛ شمارهٔ صفحه: ۱؛ \lr{YASHA}.
\subsection*{اپیدمیولوژی \lr{AR}}
\begin{itemize}
\item \textbf{\ARrange{20}{30}\% کودکان} را درگیر می‌کند.
\item احتقان، رینوره، خارش مداوم + عطسه و کنژکتیویت.
\item بیماری‌های همراه: کنژکتیویت، سینوزیت، \lr{OM}، اوتیت سروز، هایپرتروفی تونسیل و آدنوئید، اگزما.
\item افزایش ۳ برابری ریسک آسم بزرگسالی.
\item معمولاً \textbf{پیش از ۶ سالگی} شروع و تشخیص داده شده است.
\end{itemize}
\subsection*{ریسک‌فاکتورهای \lr{AR}}
\begin{itemize}
\item سابقهٔ خانوادگی آتوپی، سطح \lr{IgE} \lr{\(>100\,\mathrm{U/mL}\)} قبل از ۶ سالگی.
\item مواجههٔ جنین با دود دخانیات، سیگار قبل از ۱ سالگی.
\item مواجههٔ شدید با آلرژن‌های \lr{indoor}.
\item فرزندِ کسی که پدر و مادر سابقهٔ آتوپی دارند. \ARunclear{واژهٔ آغاز خط: فرزندِ}.
\end{itemize}
یک دورهٔ در شیرخوارگی وجود دارد که مواجههٔ زیاد با آلرژن‌ها می‌تواند باعث ایجاد آلرژی در فرد مستعد از نظر ژنتیکی شود.

سگ و گربه در سال‌های اولیهٔ زندگی از ریسک آتوپی کم می‌کنند.

طولانی شدن دورهٔ شیردهی از آتوپی پیشگیری می‌کند.

\ARpage{۲}
\textbf{حاشیهٔ اضافی:} شمارهٔ ۳ از برگهٔ زیرین در گوشهٔ پایین چپ دیده می‌شود.
\subsection*{اتیولوژی و تقسیم‌بندی‌ها}
\begin{itemize}
\item قدیم: \lr{AR} \lr{\(\longrightarrow\)} \lr{Seasonal}، \lr{Perennial}.
\item جدید: \lr{AR} \lr{\(\longrightarrow\)} \textbf{\lr{Intermittent}}، \textbf{\lr{Persistent}}.
\end{itemize}
\subsection*{شدت \lr{AR}}
\begin{latin}
\begin{itemize}
\item mild intermittent
\item moderate--severe intermittent
\item mild persistent
\item moderate--severe persistent
\end{itemize}
\end{latin}
\begin{itemize}
\item \lr{Intermittent}: \textbf{کمتر از ۴ روز در هفته یا کمتر از ۴ هفتهٔ متوالی}.
\item \lr{mild}: خواب \lr{N/L} است، فعالیت روزانه مختل نیست. مدرسه و کار را می‌رود.
\end{itemize}
\textbf{علائم \lr{Intermittent AR} معمولاً با گذشت زمان رفع می‌شوند.}

\lr{AR}های مکرر منجر به فرایند \lr{priming} می‌شوند، یعنی موارد بعدی شدیدتر خواهند بود.
\subsection*{تظاهرات بالینی \lr{AR}}
\begin{itemize}
\item \lr{Blow nose} در کودکان بزرگ‌تر؛ شیرخواران: \lr{Snort}، \lr{Sniff}.
\item خارش بینی \lr{\(\longrightarrow\)} دستکاری \lr{\(\longrightarrow\)} اپیستاکسی.
\item \lr{Allergic Salute}.
\end{itemize}

\ARpage{۳}
\begin{itemize}
\item احتقان بینی، عطسه، رینورهٔ شفاف.
\item \textbf{مواجههٔ شدیدتر \lr{\(\longrightarrow\)} علائم شدیدتر}.
\item از دست دادن موقت بویایی.
\item سردرد، ویزینگ، سرفه.
\end{itemize}
چون احتقان در خواب تشدید می‌شود \lr{\(\longrightarrow\)} تنفس دهانی \lr{\(\longrightarrow\)} اختلال عملکردی.
\begin{latin}
\begin{itemize}
\item Abnormal facial development
\item Dental malocclusion
\item Allergic gape \(\longrightarrow\) Continuous open mouth breathing
\item Chapped lips
\item Allergic shiners (Dark circle under the eyes)
\end{itemize}
\end{latin}
\subsection*{تشخیص‌های افتراقی}
\begin{itemize}
\item افتراق \lr{AR} اینترمیتنت و پرسیستنت: \textbf{شرح حال و \lr{SPT}}.
\item \lr{Nonallergic rhinitides}: علائم اسپورادیک دارند.
\item رینیت غیرآلرژیک التهابی: علائم و درمان مشابه \lr{AR} ولی بدون افزایش سطح \lr{IgE}.
\item \textbf{رینیت وازوموتور}: به دنبال فعالیت فیزیکی.
\item رینیت عفونی.
\item اختلالات ساختاری: پولیپ بینی، انحراف سپتوم.
\item رینیت \lr{medicamentosa}: استفادهٔ زیاد از وازوکانستریکتورهای موضعی.
\item رینیت هورمونال: بارداری، هایپوتیروئیدی. [در ادامه دو بخش خط‌خورده و ناخوانا وجود دارد.]
\item واسکولیت‌ها.
\item اختلالات گرانولوماتوز.
\end{itemize}

\ARpage{۴}
\textbf{حاشیهٔ اضافی:} شمارهٔ ۵ و تکه‌هایی از نوشته‌های برگهٔ زیرین، از جمله \lr{Su}، در سمت چپ دیده می‌شود؛ بخش‌های دیگر پوشیده‌اند.
\subsection*{عوارض \lr{AR}}
\begin{itemize}
\item \textbf{عدم درمان صحیح آن آسم زمینه‌ای را تشدید می‌کند.}
\item بیشترین همراهی با کنژکتیویت آلرژیک: در ۷۰\% \lr{AR}ها.
\item خارش، قرمزی، التهاب ملتحمه.
\item سینوزیت مزمن یک عارضهٔ شایع آن است.
\item گاهی همراه با ترشحات چرکی.
\item \textbf{\lr{AERD}} (بیماری تنفسی تشدید شده به دنبال آسپرین).
\item رینوسینوزیت مزمن + پولیپ بینی، آسم، حساسیت به آسپرین.
\item اغلب پاسخ به درمان آن ضعیف است.
\item ۷۸\% بیماران آسم، \lr{AR} دارند.
\item ۳۸\% بیماران \lr{AR}، آسم دارند.
\item \lr{PND} در زمینهٔ \lr{AR} سرفهٔ مزمن ایجاد می‌کند.
\item \lr{OME} عارضهٔ شایع \lr{AR} است.
\end{itemize}
\lr{AR} در برخی افراد طی نوجوانی کاهش می‌یابد ولی در ابتدا جوانی دوباره برمی‌گردد تا اینکه \textbf{در دههٔ ۵ام زندگی رفع می‌شود}.
\subsection*{ارزیابی کیفیت زندگی برای شروع درمان}
\begin{itemize}
\item در \ARrange{6}{12} سال: \lr{PRQLQ}.
\item در \ARrange{12}{17} سال: \lr{ARQLQ}.
\end{itemize}

\ARpage{۵}
\subsection*{یافته‌های آزمایشگاهی}
\begin{itemize}
\item \textbf{\lr{SPT} بهترین تست برای شناسایی \lr{sIgE}}.
\item حساسیت به آلرژن‌های تنفسی فصلی نیاز به حداقل ۲ سال مواجهه دارد.
\item \textbf{مونته‌لوکاست از ۱ روز قبل قطع شود}.
\item آنتی‌هیستامین نسل یک از \ARrange{3}{6} روز و نسل دو از \ARrange{5}{7} روز [قبل قطع شود]. \ARunclear{رقم دوم بازهٔ نسل یک: ۶ یا ۴}.
\item اگر بیمار درماتیت منتشر یا درماتوگرافیسم داشت: تست سرمی \lr{sIgE}؛ یا ریسک آنافیلاکسی.
\item اسمیر ترشحات بینی: وجود ائوزینوفیل کمک به تشخیص.
\item وجود نوتروفیل: به نفع رینیت عفونی.
\item ائوزینوفیلی، سطح \lr{total IgE} ارزش تشخیصی پایینی دارند.
\end{itemize}
\subsection*{درمان \lr{AR}}
\begin{itemize}
\item قبل از شروع درمان \lr{QoL} را ارزیابی کنید.
\item \textbf{کاهش مواجهه با آلرژن‌های \lr{indoor}}.
\item کاور ضدآلرژن بالش و تخت.
\item شستشوی \textbf{هفتگی} با آب ۵۵ درجهٔ بالش و پتو.
\item حذف حیوانات خانگی.
\item فیلتر هوای \lr{HEPA}.
\item عطسه، رینوره، علائم چشمی: آنتی‌هیستامین خوراکی.
\item تجویز \lr{PRN}.
\item بیماری خفیف ـ متوسط را به تنهایی کنترل می‌کند.
\end{itemize}

\ARpage{۶}
\begin{itemize}
\item \textbf{خط اول: آنتی‌هیستامین نسل ۲، چون کمتر خواب‌آور}.
\item برخی ترکیب‌ها سودوافدرین همراه دارند.
\item رفع احتقان بینی و سینوس.
\item خط اول نیست و باید با احتیاط استفاده شود.
\end{itemize}
\begin{itemize}
\item آنتی‌هیستامین نسل ۱: بروم‌فنیرامین، کلرفنیرامین، تریپرولیدین.
\item آنتی‌هیستامین نسل ۲: دس‌لوراتادین، فکسوفنادین.
\item آنتی‌پایرتیک: استامینوفن، ایبوپروفن.
\item ضدسرفه (آنتی‌توسیو): \lr{guaifenesin}، دکسترومتورفان.
\item آنتی‌کولینرژیک: \lr{Methscopolamin}.
\end{itemize}
\subsection*{فرم خوراکی سودوافدرین}
\begin{itemize}
\item تحریک‌پذیری، بی‌خوابی.
\item مرگ شیرخوار.
\item \textbf{زیر ۴ سال ممنوع}.
\end{itemize}
\subsection*{آنتاگونیست رسپتور لکوترین خوراکی}
\begin{itemize}
\item \textbf{خط اول \lr{AR} نیست}.
\item \textbf{عوارض نوروسایکیاتریک مونته‌لوکاست}.
\end{itemize}
\subsection*{اسپری نازال ایپراتروپیوم بروماید}
\begin{itemize}
\item آنتی‌کولینرژیک.
\item در درمان \lr{Serous rhinorrhea}.
\end{itemize}
\subsection*{اسپری نازال دکونژستان}
\begin{itemize}
\item اکسیمتازولین، فنیل‌افرین.
\item کمتر از ۵ روز می‌توان استفاده کرد. حداکثر یک‌بار در ماه.
\end{itemize}

\ARpage{۷}
\begin{itemize}
\item سدیم کروموگلیکات: ضداحتقان.
\item نیازمند استفادهٔ مکرر در ۴ ساعت.
\end{itemize}
داروهای لکوترین مودیفایر: اثر مختصر بر رینوره و انسداد بینی.

\textbf{شستشو و \lr{irrigation} بینی با سالین}: مکمل درمانی خوب برای \lr{AR}.

اگر علائم شدید بود یا با داروهای بالا کنترل نشد:
\begin{itemize}
\item \lr{CS} اینترانازال؛ برای کنژکتیویت هم مؤثر است.
\end{itemize}
\lr{CS} در آن دسته رینوره‌ای مؤثر است که التهاب ائوزینوفیلیک باشد.
\begin{itemize}
\item بکلومتازون، تریامسینولون، فلونیزولاید: جذب گوارشی و تنفسی.
\item بودزوناید، فلوتیکازون، مومتازون، سیکلزوناید.
\item عملکرد موضعی بهتر و جذب سیستمیک کمتر.
\end{itemize}
\subsection*{شدیدترین فرم \lr{AR}}
\begin{itemize}
\item همزمان \lr{CS} اینترانازال و \lr{AH} خوراکی.
\end{itemize}
\subsection*{ایمونوتراپی}
\begin{itemize}
\item در کودکانی که آلرژن عامل \lr{AR} شناسایی شده ولی قابل اجتناب از مواجهه نیست؛ آلرژن باید \lr{IgE mediated} [باشد].
\item و به دارو هم جواب نداده.
\item \textbf{از بروز آسم جلوگیری می‌کند}.
\end{itemize}

\ARpage{۸}
\textbf{حاشیهٔ اضافی:} شمارهٔ ۹ از برگهٔ زیرین در پایین چپ و تکهٔ \lr{Su} در کنار چپ دیده می‌شود.
\begin{itemize}
\item \lr{Omalizumab}: آنتی‌بادی ضد \lr{IgE}.
\item آسم که کنترل شود، \lr{AR} همراه را هم کنترل می‌کند.
\item \lr{Dupilumab}: آنتاگونیست \lr{IL-4}.
\item \lr{FDA} برای رینوسینوزیت با نازال پولیپ در بزرگسالان (\lr{AERD}).
\end{itemize}
درمان \lr{AR} معمولاً کنژکتیویت آلرژیک را هم کنترل می‌کند.
\begin{itemize}
\item \lr{CS} اوکولار: قوی‌ترین درمان برای آلرژی چشمی.
\item عارضه: کاتاراکت، افزایش فشار چشم.
\item تنها زمانی که به درمان‌های قبلی جواب نداده.
\item مشاورهٔ چشم‌پزشکی.
\end{itemize}
\subsection*{پروگنوز}
\begin{itemize}
\item \ARrange{10}{23}\% کودکان وارد ریمیشن به دنبال درمان می‌شوند.
\end{itemize}
\subsection*{علل غیرآلرژیک رینیت}
\begin{itemize}
\item علل ساختاری: \lr{CSF leak}، تودهٔ نازال، جسم خارجی، هایپرتروفی آدنوئید، انحراف سپتوم.
\item علل عفونی: حاد، مزمن، سیفلیس مادرزادی.
\item دارو: \lr{OCP}، داروی ضد فشار خون، \textbf{\lr{NSAID}، آسپرین، کوکائین}.
\end{itemize}

\ARpage{۹}
\subsection*{دس‌لوراتادین}
\begin{itemize}
\item ۶ ماه تا ۱ سال: \lr{1 mg} روزانه.
\item ۱ سال تا ۵ سال: \lr{1.25 mg} روزانه.
\item ۶ سال تا ۱۱ سال: \lr{2.5 mg} روزانه.
\item بالای ۱۲ سال: \lr{5 mg} روزانه.
\end{itemize}
\subsection*{مونته‌لوکاست}
\begin{itemize}
\item ۶ ماه تا ۵ سال: \lr{4 mg} روزانه.
\item ۶ سال تا ۱۴ سال: \lr{5 mg} روزانه.
\item بالای ۱۴ سال: \lr{10 mg} روزانه.
\end{itemize}

\clearpage\section*{پیوست: تصاویر کامل منبع}
\IfFileExists{images/pic01.jpg}
  {\def\ARimagepath{images/}}
  {\def\ARimagepath{pediatrics/ped-immunology/allergic rhinitis/images/}}
\newcommand{\ARoriginal}[1]{\clearpage\begin{center}
  \textbf{تصویر اصلی \lr{#1}}\par\smallskip
  \includegraphics[width=\textwidth,height=0.89\textheight,keepaspectratio]{\ARimagepath pic#1.jpg}
\end{center}}
\ARoriginal{01}
\ARoriginal{02}
\ARoriginal{03}
\ARoriginal{04}
\ARoriginal{05}
\ARoriginal{06}
\ARoriginal{07}
\ARoriginal{08}
\ARoriginal{09}
\endgroup
\ARfinish
